\section{Related work}
\label{sec:related}

\subsection{Segmentation on CAMUS}

CAMUS~\cite{leclerc2019camus} provides apical two- and four-chamber views for
500 patients with expert delineations of the LV endocardium, LV myocardium and
left atrium at end-diastole (ED) and end-systole (ES), together with an official
patient-level split. The original study established that a plain U-Net trained
on a sufficiently large annotated set outperforms classical methods and lands,
on average, below inter-observer variability though above intra-observer
variability.

Subsequent work has largely followed the encoder--decoder line:
nnU-Net's self-configuring pipeline~\cite{isensee2021nnunet}, transformer
hybrids such as TransUNet~\cite{chen2021transunet} and
Swin-UNet~\cite{cao2022swinunet}, and more recently state-space and
Kolmogorov--Arnold variants. Reported mean Dice on this benchmark now clusters
between $0.90$ and $0.95$. Two observations motivate the present study. First,
that band straddles the $0.899$ inter-observer figure reported by the dataset's
own authors, so differences within it are difficult to interpret as anatomical
accuracy. Second, comparisons across papers are frequently not like-for-like:
protocols differ in the split used, in whether boundary metrics are computed on
a resampled grid or at native resolution, and in the hardware on which latency
is measured.

\subsection{Anatomically constrained and boundary-aware segmentation}

Because overlap saturates, several lines of work target the boundary or the
shape prior directly. Post-hoc anatomical correction enforces topological
validity on network outputs without retraining. Contour-based and
curve-deforming formulations regress vertex positions rather than dense masks.
Boundary-weighted and distance-transform losses reweight the objective toward
the contour. We test the last of these directly (Sec.~\ref{sec:results}) and
find it does not help on this benchmark, which we attribute to the contour being
precisely where the reference labels are least reliable.

\subsection{Real-time instance segmentation}

Detector-based instance segmentation descends from Mask
R-CNN~\cite{he2017maskrcnn} and, for the real-time branch, from YOLACT's
prototype-mask formulation~\cite{bolya2019yolact}: the network predicts a small
basis of mask prototypes plus per-instance coefficients and composes them at a
fraction of the input resolution. The YOLO family~\cite{redmon2016yolo} carries
this head in its modern versions and is designed for the latency regime that
bedside echocardiography requires.

The formulation carries an assumption that does not hold for cardiac anatomy:
an instance is a single closed polygon. Nested and annular structures --- of
which the myocardium is the canonical medical example --- have no such
representation. To our knowledge this obstacle has not been examined on CAMUS,
and detector-based segmentation has not been placed on the same measurement
footing as the semantic-segmentation literature.

\subsection{Reproducibility and evaluation practice}

A recurring finding in empirical machine learning is that reported improvements
frequently fall within the variance induced by random seeds, and that
comparisons are sensitive to evaluation details that papers do not always
specify. Medical imaging inherits this problem with an additional complication:
test sets are small --- 50 patients here --- so paired significance tests over
a few hundred frames are powerful enough to certify differences that a repeat
run would erase. Our measurement of the seed floor (Sec.~\ref{sec:noise})
quantifies this for CAMUS specifically, and our audit of the reference EF
implementation (Sec.~\ref{sec:ef}) shows that an evaluation fault can survive in
widely used released code precisely because it is invisible on ground truth.
